% !PRINTED 29
\weberpage{29}
Im anderen Falle wollen wir die von einander verschiedenen der Gruppen (7) mit
\weq{8}{Q_\nu, \; Q'_\nu, \; Q''_\nu \; \ldots}
bezeichnen, und daraus eine Gruppe
\weq{9}{R = (Q_\nu, \; Q'_\nu, \; Q''_\nu \; \ldots)}
ableiten, die dadurch definirt sein soll, dass es die kleinste Gruppe ist, die jede der Gruppen (8) als Theiler enthält, und die wir als das \gesperrt{kleinste gemeinschaftliche Vielfache} der Gruppen $Q_\nu, Q'_\nu, Q''_\nu \ldots$ bezeichnen können.

Dass es eine und nur eine solche Gruppe $R$ giebt, und dass jede Gruppe, die durch $Q_\nu, Q'_\nu, Q''_\nu \ldots$ theilbar ist, auch durch $R$ theilbar sein muss, ist leicht einzusehen. Denn erstens giebt es endliche Gruppen, z. B. die Gruppe $P_\nu$, die alle Gruppen (8) als Theiler enthalten, und zweitens, wenn zwei Gruppen $R, R'$ alle diese Gruppen enthalten, so gilt dasselbe auch von dem Durchschnitt von $R$ und $R'$. Ist daher $R$ von möglichst niedrigem Grade, so muss dieser Durchschnitt mit $R$ identisch sein, und $R$ ist also ein \gesperrt{Theiler von} $R'$. Wenn $R'$ auch von möglichst niedrigem Grade ist, so ist $R'$ von $R$ nicht verschieden.

Von dieser Gruppe $R$ weisen wir nun nach:

a) dass sie ein Normaltheiler von $P_{\nu-1}$ ist,

b) dass sie eine commutative Gruppe ist.

Wenn wir dies Beides nachgewiesen haben, so sind wir am Ziele; denn dann können wir, wenn es nicht noch einen commutativen Normaltheiler \gesperrt{niedrigeren Grades} von $P_{\nu-1}$ giebt, $R$ selbst für $Q_{\nu-1}$ wählen.

Das Erste ist aber ohne Weiteres klar. Denn bedeutet $\gamma$ irgend ein Element aus $P_{\nu-1}$, so ist
\[ \gamma^{-1} R \gamma = (\gamma^{-1} Q_\nu \gamma, \; \gamma^{-1} Q'_\nu \gamma, \; \gamma^{-1} Q''_\nu \gamma \; \ldots) \]
das kleinste gemeinschaftliche Vielfache der Gruppen
\[ \gamma^{-1} Q_\nu \gamma, \; \gamma^{-1} Q'_\nu \gamma, \; \gamma^{-1} Q''_\nu \gamma \; \ldots \]
Diese Gruppen aber sind nach der Definition (7) in ihrer Gesammtheit von den Gruppen
\[ Q_\nu, \; Q'_\nu, \; Q''_\nu \; \ldots \]
nicht verschieden, und folglich ist
\[ \gamma^{-1} R \gamma = R, \]
d. h. $R$ ist Normaltheiler von $P_{\nu-1}$.
